\documentclass[10pt,a4paper]{amsart}
\usepackage{amsmath,amssymb,amsthm,graphicx}
\usepackage{verbatim}
\newenvironment{reference}{\comment}{\endcomment}
%\newenvironment{reference}{}{}
\newenvironment{slogan}{\comment}{\endcomment}
\newenvironment{history}{\comment}{\endcomment}

% For commutative diagrams we use Xy-pic
\usepackage[all]{xy}

% We use 2cell for 2-commutative diagrams.
\xyoption{2cell}
% Theorem environments.
%
\theoremstyle{plain}
\newtheorem{theorem}[subsection]{Theorem}
\newtheorem{proposition}[subsection]{Proposition}
\newtheorem{lemma}[subsection]{Lemma}

\theoremstyle{definition}
\newtheorem{definition}[subsection]{Definition}
\newtheorem{example}[subsection]{Example}
\newtheorem{exercise}[subsection]{Exercise}
\newtheorem{situation}[subsection]{Situation}

\theoremstyle{remark}
\newtheorem{remark}[subsection]{Remark}
\newtheorem{remarks}[subsection]{Remarks}

\numberwithin{equation}{subsection}

% Macros
%
\def\lim{\mathop{\mathrm{lim}}\nolimits}
\def\colim{\mathop{\mathrm{colim}}\nolimits}
\def\Spec{\mathop{\mathrm{Spec}}}
\def\Hom{\mathop{\mathrm{Hom}}\nolimits}
\def\Ext{\mathop{\mathrm{Ext}}\nolimits}
\def\SheafHom{\mathop{\mathcal{H}\!\mathit{om}}\nolimits}
\def\SheafExt{\mathop{\mathcal{E}\!\mathit{xt}}\nolimits}
\def\Sch{\mathit{Sch}}
\def\Mor{\mathop{\mathrm{Mor}}\nolimits}
\def\Ob{\mathop{\mathrm{Ob}}\nolimits}
\def\Sh{\mathop{\mathit{Sh}}\nolimits}
\def\NL{\mathop{N\!L}\nolimits}
\def\CH{\mathop{\mathrm{CH}}\nolimits}
\def\proetale{{pro\text{-}\acute{e}tale}}
\def\etale{{\acute{e}tale}}
\def\QCoh{\mathit{QCoh}}
\def\Ker{\mathop{\mathrm{Ker}}}
\def\Im{\mathop{\mathrm{Im}}}
\def\Coker{\mathop{\mathrm{Coker}}}
\def\Coim{\mathop{\mathrm{Coim}}}

% Boxtimes
%
\DeclareMathSymbol{\boxtimes}{\mathbin}{AMSa}{"02}

%
% Macros for moduli stacks/spaces
%
\def\QCohstack{\mathcal{QC}\!\mathit{oh}}
\def\Cohstack{\mathcal{C}\!\mathit{oh}}
\def\Spacesstack{\mathcal{S}\!\mathit{paces}}
\def\Quotfunctor{\mathrm{Quot}}
\def\Hilbfunctor{\mathrm{Hilb}}
\def\Curvesstack{\mathcal{C}\!\mathit{urves}}
\def\Polarizedstack{\mathcal{P}\!\mathit{olarized}}
\def\Complexesstack{\mathcal{C}\!\mathit{omplexes}}
% \Pic is the operator that assigns to X its picard group, usage \Pic(X)
% \Picardstack_{X/B} denotes the Picard stack of X over B
% \Picardfunctor_{X/B} denotes the Picard functor of X over B
\def\Pic{\mathop{\mathrm{Pic}}\nolimits}
\def\Picardstack{\mathcal{P}\!\mathit{ic}}
\def\Picardfunctor{\mathrm{Pic}}
\def\Deformationcategory{\mathcal{D}\!\mathit{ef}}
\usepackage{cleveref}
\crefname{originalreference}{original reference}{original references}
\Crefname{originalreference}{Original reference}{Original references}
\makeatletter
\newlabel{lemma-good-quotient}{{09KN}{source}{}{sourcepointer-1}{}}
\newlabel{lemma-good-quotient@cref}{{[originalreference][0][]09KN}{[1][0][]source}{}{}{}}
\makeatother
\crefname{originalreference}{original reference}{original references}
\Crefname{originalreference}{Original reference}{Original references}
\makeatletter
\newlabel{lemma-hom-into-shift-dual-free}{{09K4}{source}{}{sourcepointer-2}{}}
\newlabel{lemma-hom-into-shift-dual-free@cref}{{[originalreference][0][]09K4}{[1][0][]source}{}{}{}}
\makeatother
\crefname{originalsection}{original section}{original sections}
\Crefname{originalsection}{Original section}{Original sections}
\makeatletter
\newlabel{section-modules-noncommutative-differential-graded}{{0FQD}{source}{}{sourcepointer-3}{}}
\newlabel{section-modules-noncommutative-differential-graded@cref}{{[originalsection][0][]0FQD}{[1][0][]source}{}{}{}}
\makeatother
\title{Embedding a differential graded module into a product of shifts of the dual algebra preserving cokernels of the differential}
\author{The Stacks Project Authors}
\date{stacks-project:dga.tex:local-snapshot-20261012; GNU FDL 1.2 or later}
\begin{document}
\maketitle
\noindent\textit{Editorial note.} All delivered mathematical text is the author's own, copied from the pinned Stacks Project chapter file sources/stacks-dga/pinned-original-source.tex; no statement or proof step was rewritten, reordered or removed.\par\medskip
\noindent\textit{Editorial note.} The counted proof is the author's complete argument: the free-module surjection onto M-dual (parallel to the proof of lemma-good-quotient, cited in place), the double-dual embedding into the product of the I\_k, the dualized two-term sequence showing each I\_k is an extension of shifts of A-dual, and the functional-separation argument for property (2) through lemma-hom-into-shift-dual-free are all retained verbatim.\par\medskip
\noindent\textit{Editorial note.} The label lemma-good-sub names the counted statement, and the proof environment immediately following it in the source is delivered verbatim as its proof.\par\medskip
\noindent\textit{Editorial note.} The counted statement is carried by the author's own lemma environment in the chapter source.\par\medskip
\noindent\textit{Allowed standard prerequisites (editorial).}\begin{itemize}
\item The duality functors between left and right differential graded modules. Same chapter, Section section-modules-noncommutative-differential-graded. Supplies the duality N maps to N-dual with all its properties, through which the embedding into a product of duals is built.
\item Free surjections onto duals and Hom into shifted duals. Same chapter, Lemma lemma-good-quotient (its proof) and Lemma lemma-hom-into-shift-dual-free. The parallel free-module construction produces the surjection onto M-dual, and the Hom-into-shift-dual lemma turns a separating functional on the cokernel into a map to a shift of A-dual, giving property (2).
\end{itemize}\medskip
\ifdefined\hypertarget\hypertarget{sourcepointer-1}{}\fi
\noindent\textit{Source reference.} Original reference 09KN: sources/stacks-dga/pinned-original-source.tex, source lines 2490--2490. Prerequisite pointer invoked at proof line(s) 2680; it names a standard fact of the same chapter that the counted argument uses, and every step of the counted proof built on it is retained verbatim in the delivered body.\par\medskip
\ifdefined\hypertarget\hypertarget{sourcepointer-2}{}\fi
\noindent\textit{Source reference.} Original reference 09K4: sources/stacks-dga/pinned-original-source.tex, source lines 2366--2366. Prerequisite pointer invoked at proof line(s) 2700; it names a standard fact of the same chapter that the counted argument uses, and every step of the counted proof built on it is retained verbatim in the delivered body.\par\medskip
\ifdefined\hypertarget\hypertarget{sourcepointer-3}{}\fi
\noindent\textit{Source reference.} Original section 0FQD: sources/stacks-dga/pinned-original-source.tex, source lines 2306--2306. Prerequisite pointer invoked at proof line(s) 2674; it names a standard fact of the same chapter that the counted argument uses, and every step of the counted proof built on it is retained verbatim in the delivered body.\par\medskip
\begin{lemma}
\label{lemma-good-sub}
Let $(A, \text{d})$ be a differential graded algebra.
Let $M$ be a differential graded $A$-module. There exists a homomorphism
$M \to I$ of differential graded $A$-modules with the following
properties
\begin{enumerate}
\item $M \to I$ is injective,
\item $\Coker(\text{d}_M) \to \Coker(\text{d}_I)$ is injective,
and
\item $I$ sits in an admissible short exact sequence
$0 \to I' \to I \to I'' \to 0$ where $I'$, $I''$ are products
of shifts of $A^\vee$.
\end{enumerate}
\end{lemma}

\begin{proof}
We will use the functors $N \mapsto N^\vee$ (from left to right
differential graded modules and from right to left differential
graded modules) constructed in
Section \ref{section-modules-noncommutative-differential-graded}
and all of their properties.
For every $k \in \mathbf{Z}$ let $Q_k$ be the free left $A$-module with
generators $x, y$ in degrees $k$ and $k + 1$. Define the structure of a
left differential graded $A$-module on $Q_k$ by setting $\text{d}(x) = y$
and $\text{d}(y) = 0$. Arguing exactly as in the proof of
Lemma \ref{lemma-good-quotient} we find a surjection
$$
\bigoplus\nolimits_{i \in I} Q_{k_i} \longrightarrow M^\vee
$$
of left differential graded $A$-modules. Then we can consider the injection
$$
M \to (M^\vee)^\vee \to (\bigoplus\nolimits_{i \in I} Q_{k_i})^\vee =
\prod\nolimits_{i \in I} I_{k_i}
$$
where $I_k = Q_{-k}^\vee$ is the ``dual'' right differential graded $A$-module.
Further, the short exact sequence $0 \to A[-k - 1] \to Q_k \to A[-k] \to 0$
produces a short exact sequence
$0 \to A^\vee[k] \to I_k \to A^\vee[k + 1] \to 0$.

\medskip\noindent
The result of the previous paragraph produces $M \to I$
having properties (1) and (3). To obtain property (2), suppose
$\overline{m} \in \Coker(\text{d}_M)$ is a nonzero element of
degree $k$. Pick a map $\lambda : M^k \to \mathbf{Q}/\mathbf{Z}$
which vanishes on $\Im(M^{k - 1} \to M^k)$ but not on $m$. By
Lemma \ref{lemma-hom-into-shift-dual-free}
this corresponds to a homomorphism $M \to A^\vee[k]$ of
differential graded $A$-modules which does not vanish on $m$.
Hence we can achieve (2) by adding
a product of copies of shifts of $A^\vee$.
\end{proof}

\end{document}
